% !TeX root = ../Masterarbeit.tex

\documentclass[Masterarbeit.tex]{subfile}

\begin{document}
	
	\chapter{Introduction}
	
	\section{Motivation and Problem}
	
	Predicting changes of Earth system variables on interannual to decadal scales has been the focus of large amounts of current climate research. From forecasts of extreme weather events like droughts or heat waves to predictions of precipitation and temperature on regional scales, decadal predictions are the center of a lot of scientific debate \citep{marotzke_miklip_2016, keenlyside_prospects_2010}. Land surface and atmospheric components of the Earth system often change so quickly that they rarely possess any memory beyond seasonal changes. Thus, predictions on larger timescales from months to decades cannot solely rely on them as indicators for future changes. Contrary to its land and atmospheric components, the global oceans encompass low frequency variability and  a much longer memory. Therefore, they play an important role in climate modeling on interannual to decadal time scales. In particular ocean heat content (OHC) indicates long term changes in heat fluxes between the Earth system's oceanic and atmospheric components \citep{cheng_improved_2017}. In addition to that, OHC also shows a consistent and slowly varying response to changes in external boundary conditions of the climate system. Therefore, representing the overall changes in ocean heat content as well as its variability and response to external forcings is one of the major challenges for decadal climate predictions.\\
	Ocean heat content is not only important for climate predictions due to its low frequency variability. In comparison to surface land and sea temperatures, changes in ocean heat content also show lower background noise in comparison to their signal, making them a better indicator for climatic changes \citep{cheng_decadal_2018}. Additionally, OHC is instrumental to the Earth's energy budget and its response to anthropogenic greenhouse gas emissions. Greenhouse gases reflect longwave radiation back towards the Earth's surface, creating an imbalance in the Earth's energy budget in the form of a surplus of incoming energy compared to the outgoing longwave radiation. More than $90 \%$ of this energy imbalance is retained in the global oceans with over $30 \%$ being stored in the North Atlantic \citep{hakkinen_heat_2015}. This does not only lead to thermal expansion and thus a direct rise in sea level but also affects the formation of sea ice and the global circulation systems. Contrary to most of the global oceans, not all of the North Atlantic OHC is rising. Especially in the subpolar gyre (SPG), OHC has been overall declining over the past 70 years. Most of this change can be attributed to natural variability related to the North Atlantic Oscillation (NAO) \citep{lozier_spatial_2008}. This natural variability in the North Atlantic and its connection to Norther European and American regional climate is another reason, why the study of North Atlantic and SPG OHC has gotten a lot of attention in recent literature.\\
	The North Atlantic and especially its sea surface temperature (SST) can be predicted with relatively high skill on interannual to decadal timescales \citep{buckley_predictability_2019}. However, especially when examining subsurface temperatures and OHC in the SPG, larger uncertainties remain. Better estimating these uncertainties could thus give insights into the different regional OHC trends in the North Atlantic and their impact on European climate and weather. \\
	In general, ocean heat content estimates are dependent on several factors. A first and very direct source of uncertainties is the bias measurement instruments bring into the calculation of ocean heat content. Over large parts of the $20^{\un{th}}$ century, mechanical bathythermographs (MBTs) and expendable bathythermographs (XBTs) provided considerable amounts of ocean temperature profile data. MBTs accounted for $68 \%$ of profiles from 1941--1966 before being replaced by XBTs, which then accounted for $38 \%$ of profile data from 1970--2001 \citep{gouretski_correction_2020}. Both measurement methods have been found to incorporate a time variant warm bias, which significantly impacted ocean heat content timeseries. \cite{gouretski_correction_2020} show that time and spatially dependent corrections can help significantly reduce the MBT bias of $0.2\degree$C. Nonetheless, even after the correction a bias of up to $0.1\degree$C remains, particularly in the southern hemisphere. Similarly, \cite{cheng_time_2014} have designed a correction method for the XBT bias, which limits the bias up to $\pm 0.04\degree$C. In addition to XBT and MBT biases, \cite{boyer_sensitivity_2016} have identified mapping methods and baseline climatologies as additional sources of uncertainties. Especially since scientific consensus on the choice of XBT bias correction method has been reached by now \citep{garcia-soto_overview_2021}, \cite{boyer_sensitivity_2016} have identified the choice of mapping method the largest source of uncertainty for the calculation of ocean heat content. \\
	
	Mapping methods refer to the way of reconstructing a full set of values of a variable from a limited amount of observations. As ocean heat content cannot be observed directly, it is usually calculated by integrating over temperature profiles from the ocean surface down to a depth of up to 2.000 m. Below 2.000 m, very few observations exist, as measurement instruments, such as Argo floats, do not descend below this threshold \cite{wong_argo_2020}. Thus, even sophisticated mapping methods have significant difficulties in estimating deeper ocean heat content. However, even above 2.000 m depth, subsurface temperature observations often lack adequate coverage and thus cannot be used without further processing to initialize climate models. In order to create a complete set of values close to the observed state, a range of statistical mapping methods can be applied. On the one hand, statistical methods interpolate between missing data points on a purely mathematical basis. Statistical interpolation approaches, such as kriging, are relatively cost effective and have shown reasonable success in reconstruction temperature data \citep{varentsov_high-resolution_2020}. However, most statistical infilling has shown difficulties to map variables in undersampled regions. This difficulty mostly stems from their relaxation back to the underlying baseline climatology without considering larger dynamical trends or physical processes. The other larger field of mapping methods, data assimilation, tries to counter these difficulties by aiming to bring together Earth system model states and real life observations. It assumes both to be inaccurate and incomplete representations of the state of a system. Data assimilation then tries to incorporate both into a new statistical estimation of the true state. Of the many data assimilation methods, this study will focus on Ensemble Kalman Filter (EnKF) assimilations. EnKFs provide computationally efficient and very accurate results compared to other assimilation methods, especially in high dimensional systems, such as the climate and Earth system \citep{carrassi_data_2018}. A detailed overview of the mechanisms behind EnKFs as well as other assimilation methods can be found in \cite{carrassi_data_2018}.
	
	\begin{figure}[h]
		\centering
		\includegraphics[width=0.8\textwidth]{plots/ohc_timeseries.png}
		\caption{Estimates of upper 2.000 m North Atlantic ocean heat content using different data assimilation and interpolation methods \citep{hakkinen_heat_2015}.}
		\label{fig:ohc_timeseries}
	\end{figure}
	
	In general, statistical interpolation and data assimilations have shown high correlation to each other in the reconstruction of oceanic heat content \cite{brune_time_2018}. However, in the $20^{\text{th}}$ century, data assimilation runs show lower correlation to interpolation methods compared to the last 20 years \citep{brune_time_2018}. The biggest difference between both time periods for the mapping of OHC is the amount of subsurface temperature observations available in the reconstruction process. At the start of the 2000s the amount of subsurface temperature observations in global oceans started to rise rapidly due to the introduction of the Argo project. The Argo project is an internationl cooperation, which placed 4.000 swimming buoys in the global oceans \citep{roemmich_Argo_2009}. These buoys drift with ocean currents and continually measure temperatures up to a depth of 2.000 m. Due to their significant impact on the availability of NA subsurface temperature observations, the time before the introduction of Argo floats in 2004 will from hereon out be referred to as pre-Argo era, while the time from 2004 onwards will be referred to as Argo era. The timeseries of different infilling methods of global upper ocean heat content show much higher correlation with each other during the Argo era, as can be seen in figure \ref{fig:ohc_timeseries}. As there is no reference point for the ground truth available at points without direct observation, the only estimate of the missing data comes in form of uncertainty ranges of different infilling methods. Therefore, the uncertainty range of ocean heat content estimates at locations with undersampled temperature measurements is much during the pre-Argo era \citep{levitus_world_2012}. \\
	
	
	\begin{figure}
		\centering
		\includegraphics[width=1.0\linewidth]{plots/Ensemble_spread}
		\caption{Ensemble spread of an Ensemble Kalman Filter (EnKF) data assimilation of North Atlantic OHC.}
		\label{fig:ensemblespread}
	\end{figure}

	This uncertainty is also visible when examining one specific assimilation ensemble. Figure \ref{fig:ensemblespread} shows the OHC estimate of a 16 member EnKF assimilation ensemble. The technical details of this assimilation are described in detail in section \ref{assimilation}. The ensemble spread, which represents the uncertainty within the ensemble, is significantly larger during the pre-Argo era than during the 2000s. But assimilation ensembles do not only show a higher overall spread in times with lower observational density. Missing observations can also lead to the assimilation not being able to compensate for regional model uncertainties adn errors through the integration of observations \cite{vimal_personal_2022}. In addition to rising uncertainty with fewer observations, EnKF assimilations also show several other limitations. EnKFs require high computational resources to assimilate large datasets with high resolutions. Additionally, they are unflexible when changing an overall timeseries of OHC estimates. As an ensemble approach, none of the members can be calculated or changed independent from the others. Thus, EnKF assimilations require an overall recalculation of the entire ensembles when they for example are required to generate additional ensemble members. Similarly, when their input data is changed, for example by a new bias correction for in situ measurements or the exclusion of some other observations, EnKF assimilations need to recalculate the entire ensemble. \\
	
	In recent years, data assimilation techniques have often been combined with machine learning approaches, more specifically with neural networks, to improve the reconstruction of full system states. Many computer scientists have aimed at replacing data assimilation completely by machine learning \citep{geer_learning_2021}. However, others have aimed at combining assimilation methods with machine learning reconstructions in order to improve on both individual methods \citep{brajard_combining_2021, gottwald_combining_2021}. Machine learning has the advantage of combining numerical infilling with Earth system model data \citep{kadow_artificial_2020}. Especially convolutional neural networks have proven capable of capturing spatial patterns of their training input. These abilities led to a lot of research assessing the the theoretical potential to replace data assimilation tasks by machine learning \citep{geer_learning_2021}. These machine learning methods have often been transferred from the field of image inpainting. The considerable advantage of image inpainting deep learning techniques in comparison to statistical inpainting methods is their partial awareness of image semantics through the training process \citep{liu_image_2018}. In 2020, \cite{kadow_artificial_2020} managed to transfer deep learning techniques developed for image inpainting to the field of climate sciences. They adapted the inpainting techniques developed by \cite{liu_image_2018} to reconstruct incomplete global 2m-air temperature data. The reconstructed Climate Model Intercomparison Project Phase 5 (CMIP5) data showed lower root mean squared errors (rsme) than both conventional kriging as well as PCA mapping methods. Furthermore, the convolutional U-net used by \cite{kadow_artificial_2020} was able to successfully reconstruct the historically well recorded 1877 El Nino event from observational data, which both PCA and kriging reconstructions did not reveal. Furthermore, \cite{penny_integrating_2022} succesfully used recurrent neural networks (RNNs) for data assimilation in order to initialize short-term forecasts. Similarly, \cite{arcucci_deep_2021} suggest that training an RNN with data assimilations of observational states can significantly reduce the model error in what they call deep data assimilation (DDA).\\
	
	In comparison to traditional assimilations, machine learning approaches can reduce the computational costs significantly especially when working with large datasets \citep{arcucci_deep_2021}. Furthermore, machine learning approaches can integrate changes in observational data and new ensemble members much more easily as they are independent of an entire model setup. Careful selection of training data from periods with high observational density can lastly also limit model uncertainties to a minimum, as a trained network might be able to transfer trained patterns from periods with many observations to periods with few. \\
	In order to utilize these advantages, I will build on the theoretical work of \cite{arcucci_deep_2021} by training a neural network based on the one used by \cite{kadow_artificial_2020} to reconstruct three dimensional North Atlantic temperature data. This approach aims at answering two main questions surrounding the current estimates of SPG OHC. On the one hand, it aims at assessing the potential of neural networks to reproduce and even replace EnKF data assimilation of SPG OHC. Training a neural network on EnKF assimilations, its ability to reproduce assimilation estimates of SPG OHC will be tested. On the other hand, it will be ascertained if a new neural network based SPG OHC estimate can help better estimate the uncertainty of EnKF assimilations during the pre-Argo era. By only training the neural network with assimilation data from the Argo era, it can learn the more accurate assimilation patterns from times with high observational density. Transferring these to earlier time periods, it is examined if the neural network can show inconsistencies in the assimilation caused by missing observations. Additionally training a neural network on assimilation data and then reconstruction observations can help show inconsistencies and disagreements between observations and assimilation at the grid points with observations. \\
	
	To assess, if convolutional U-nets can be considered a valid alternative to EnKF data assimilation, the neural network will be trained on data assimilations during the well observed Argo era (2004 -- 2020) in order to learn on assimilations, which posses lower uncertainty and have integrated more observational data. After the network has been fully trained, visible at the progression of its training curve, it can be given masked temperature profiles for reconstruction. The network will not be trained with all 16 ensemble members. Instead, 4 ensemble members will be held back for later evaluation. This process helps to ensure that the network does not commit to overtraining on its training data. At first, the network will reconstruct the Argo era masked assimilations from one of the held back ensemble members. Each month of assimilation data will be masked with the corresponding observational pattern from that time. Given the high observational density during the Argo era, this reconstruction can be considered the first step for the neural network towards the reproduction of the assimilation OHC estimate. This approach can then be extended to the pre-Argo era, if the neural network is able to reproduce the assimilation ensemble within its uncertainty range during the Argo era. Analyzing the reconstructed OHC timeseries as well as maps from individual timesteps, the machine learning reconstructions can also be used to discuss inconsistencies and uncertainties within the assimilation. Differences between both OHC estimates can then be taken as indicator for uncertainties that are not mirrored within the assimilation's ensemble spread. Finally, I will examine, whether the neural network is able to directly reconstruct observations in the pre-Argo and Argo era as well as more recent observations it has not trained on within the assimilations uncertainty range. The results will then be evaluated in regard to the neural network's potential to reproduce or replace EnKF assimilations. Additionally, its implications on a new uncertainty estimate on the EnKF assimilations OHC outside of its ensemble spread will be considered. \\
	
	In summary, this thesis will focus on examining whether neural networks provide a valid alternative to EnKF assimilations when mapping out SPG OHC, while at the same time trying to better estimate their uncertainties. Special attention will be given to the analysis of the assimilation uncertainty during the pre-Argo era. \\
	In order to reach these aims, at first an overview of the involved data and methods will be given in chapter \ref{data_methods}. Both the EnKF assimilation used as training data and the observations utilized in these assimilations and for the later reconstructions will be described in detail in section \ref{data}. Additionally, two other OHC estimates, those of the  and the Institute for Atmospheric Physic's (IAP), will be reviewed to later provide some context to the reconstruction results. Afterwards, the utilized neural network will be described in section \ref{methods}. Then, the results of the neural network reconstructions will be described, analyzed and interpreted in chapter \ref{results}. Finally, these results will be discussed in regard to their limitations and the underlying assumptions in chapter \ref{discussion} before drawing final conclusions in chapter \ref{conclusions}. \\
	
\end{document}